\documentclass[10pt,a4paper]{article}

\usepackage[utf8]{inputenc}
\usepackage[T1]{fontenc}
\usepackage[french]{babel}
\usepackage[left=1.5cm,right=1.5cm,top=1.5cm,bottom=1.5cm]{geometry}
\usepackage{multirow,tabularx,booktabs,array}
\usepackage[table]{xcolor}
\usepackage{fouriernc}
\usepackage{amsmath}
\usepackage{fancyhdr}
\usepackage{colortbl}

\setlength{\extrarowheight}{1mm}
\renewcommand{\arraystretch}{1.35}
\setlength{\parindent}{0pt}

\pagestyle{fancy}
\renewcommand\headrulewidth{1pt}
\fancyhead[L]{Lycée Denis Diderot}
\fancyhead[R]{Terminale Bac Pro -- TRPM}
\fancyfoot[C]{}
\fancyfoot[R]{Année 2025-2026}
\fancyfoot[L]{Alexis RUIZ}

% Couleurs
\definecolor{exoheader}{RGB}{30,90,160}
\definecolor{exolight}{RGB}{220,235,255}
\definecolor{totalrow}{RGB}{255,230,200}
\definecolor{grandtotal}{RGB}{50,100,160}

% Macro en-tête de partie pleine largeur
\newcommand{\exotitle}[2]{%
  \noindent\colorbox{exoheader}{%
    \parbox{\dimexpr\linewidth-2\fboxsep\relax}{%
      \textcolor{white}{\textbf{\large #1 \hfill #2}}%
    }%
  }%
}

\begin{document}

\begin{center}
  {\LARGE\textbf{GRILLE DE CORRECTION}}\\[1mm]
  {\large Propagation du son -- Activité documentaire \hspace{1cm} Terminale Bac Pro TRPM \hspace{1cm} \textbf{Total : /20 points}}
\end{center}

\vspace{3mm}

% ===================== PARTIE A =====================
\exotitle{Partie A -- Nature et propriétés du son}{/5 pts}

\vspace{1mm}

\begin{tabularx}{\linewidth}{|c|X|c|}
  \hline
  \rowcolor{exolight}
  \textbf{Q} & \textbf{Éléments de réponse attendus} & \textbf{Pts} \\
  \hline
  1
  & Le son ne se propage pas dans le vide (onde mécanique, milieu matériel nécessaire). Citation du document.
  & 1 \\
  \hline
  \multirow{2}{*}{2}
  & Gaz : faible ($\approx$~340~m/s) / Liquides : moyenne ($\approx$~1~480~m/s) / Solides : élevée (jusqu'à 6~300~m/s)
  & 1 \\
  \cline{2-3}
  & Conclusion : le son se propage le plus vite dans les \textbf{solides}
  & 1 \\
  \hline
  \multirow{2}{*}{3}
  & Formule de conversion correcte : $v = 340 \times 3{,}6$
  & 0{,}5 \\
  \cline{2-3}
  & \textbf{Résultat :} $v = 1\,224$~km/h (avec unité)
  & 1{,}5 \\
  \hline
  \rowcolor{totalrow}
  \multicolumn{2}{|r|}{\textbf{Sous-total Partie A}} & \textbf{/5} \\
  \hline
\end{tabularx}

\vspace{4mm}

% ===================== PARTIE B =====================
\exotitle{Partie B -- Célérité dans les matériaux et CND}{/6 pts}

\vspace{1mm}

\begin{tabularx}{\linewidth}{|c|X|c|}
  \hline
  \rowcolor{exolight}
  \textbf{Q} & \textbf{Éléments de réponse attendus} & \textbf{Pts} \\
  \hline
  \multirow{2}{*}{4}
  & Relevé correct : Acier 5~900~m/s / Aluminium 6~300~m/s
  & 1 \\
  \cline{2-3}
  & Comparaison : célérité plus élevée dans l'aluminium
  & 1 \\
  \hline
  \multirow{2}{*}{5}
  & Définition : on ne détruit/n'endommage pas la pièce lors du contrôle
  & 1 \\
  \cline{2-3}
  & Intérêt TRPM : contrôle de 100\% de la production, pièces réutilisables si conformes
  & 1 \\
  \hline
  6
  & Trajet aller-retour expliqué : $\Delta t$ = durée pour 2 fois la distance, donc on divise par 2
  & 2 \\
  \hline
  \rowcolor{totalrow}
  \multicolumn{2}{|r|}{\textbf{Sous-total Partie B}} & \textbf{/6} \\
  \hline
\end{tabularx}

\vspace{4mm}

% ===================== PARTIE C =====================
\exotitle{Partie C -- Application : contrôle d'une pièce en acier}{/9 pts}

\vspace{1mm}

\begin{tabularx}{\linewidth}{|c|X|c|}
  \hline
  \rowcolor{exolight}
  \textbf{Q} & \textbf{Éléments de réponse attendus} & \textbf{Pts} \\
  \hline
  \multirow{4}{*}{7}
  & Conversion correcte : $\Delta t_B = 20{,}3 \times 10^{-6}$~s
  & 0{,}5 \\
  \cline{2-3}
  & Application de la formule : $d = \dfrac{v \times \Delta t_B}{2}$
  & 1 \\
  \cline{2-3}
  & Calcul numérique correct : $d = \dfrac{5\,900 \times 20{,}3 \times 10^{-6}}{2}$
  & 1 \\
  \cline{2-3}
  & \textbf{Résultat :} $d \approx 59{,}9$~mm -- Comparaison avec 60~mm (cohérent, écart $<$ 0,2\%)
  & 1{,}5 \\
  \hline
  \multirow{3}{*}{8}
  & Conversion correcte : $\Delta t_A = 8{,}1 \times 10^{-6}$~s
  & 0{,}5 \\
  \cline{2-3}
  & Application de la formule avec $\Delta t_A$
  & 1 \\
  \cline{2-3}
  & \textbf{Résultat :} $d \approx 23{,}9$~mm (avec unité)
  & 1{,}5 \\
  \hline
  \multirow{2}{*}{9}
  & Identification de l'anomalie interne à 23,9~mm (défaut détecté)
  & 1 \\
  \cline{2-3}
  & Justification du rejet : risque de rupture / non-conformité au cahier des charges
  & 1 \\
  \hline
  \rowcolor{totalrow}
  \multicolumn{2}{|r|}{\textbf{Sous-total Partie C}} & \textbf{/9} \\
  \hline
\end{tabularx}

\vspace{6mm}

% ===================== TOTAL =====================
\noindent\colorbox{grandtotal}{%
  \parbox{\dimexpr\linewidth-2\fboxsep\relax}{%
    \textcolor{white}{\textbf{\Large \hfill TOTAL \hfill /20 points \hfill}}%
  }%
}

\vspace{6mm}

% ===================== TABLEAU ÉLÈVES =====================
\noindent\textbf{Récapitulatif de la classe :}

\vspace{2mm}

\begin{tabularx}{\linewidth}{|>{\centering\arraybackslash}p{4.5cm}|>{\centering\arraybackslash}X|>{\centering\arraybackslash}p{1.8cm}|>{\centering\arraybackslash}p{1.8cm}|>{\centering\arraybackslash}p{1.8cm}|>{\centering\arraybackslash}p{1.8cm}|}
  \hline
  \rowcolor{exolight}
  \textbf{Nom Prénom} & \textbf{Observations} & \textbf{Part. A /5} & \textbf{Part. B /6} & \textbf{Part. C /9} & \textbf{Total /20} \\
  \hline
  & & & & & \\[6mm]
  \hline
  & & & & & \\[6mm]
  \hline
  & & & & & \\[6mm]
  \hline
  & & & & & \\[6mm]
  \hline
  & & & & & \\[6mm]
  \hline
  & & & & & \\[6mm]
  \hline
  & & & & & \\[6mm]
  \hline
  & & & & & \\[6mm]
  \hline
  & & & & & \\[6mm]
  \hline
  & & & & & \\[6mm]
  \hline
  & & & & & \\[6mm]
  \hline
  & & & & & \\[6mm]
  \hline
  & & & & & \\[6mm]
  \hline
  & & & & & \\[6mm]
  \hline
  & & & & & \\[6mm]
  \hline
  & & & & & \\[6mm]
  \hline
  & & & & & \\[6mm]
  \hline
  & & & & & \\[6mm]
  \hline
  & & & & & \\[6mm]
  \hline
  & & & & & \\[6mm]
  \hline
  \rowcolor{totalrow}
  \textbf{Moyenne} & & & & & \\
  \hline
\end{tabularx}

\end{document}
